%! suppress = EscapeAmpersand
%! suppress = TooLargeSection
%! suppress = MissingLabel

% Preamble — adapted to the layout of
% "Информатика и автоматизация / Informatics and Automation" (SPIIRAS, SPC RAS),
% https://ia.spcras.ru/index.php/sp/aims-and-scope
\documentclass[12pt,a4paper]{article}

% Packages
\usepackage[T2A,T1]{fontenc}
\usepackage[utf8]{inputenc}
\usepackage[english,russian]{babel}
\usepackage{amsmath}
\usepackage{amssymb}
\usepackage{amsthm}
\usepackage{xcolor}
\usepackage{multicol}
\usepackage{proof}
\usepackage{mathtools}
\usepackage{hyperref}
\hypersetup{colorlinks=true, linkcolor=black, citecolor=black, urlcolor=blue}

\usepackage{geometry}
% Journal page setup: A4, 2.0 cm margins, 1.5 line spacing.
\geometry{a4paper, top=20mm, bottom=20mm, left=20mm, right=20mm}

\usepackage{setspace}
\onehalfspacing

\usepackage{enumitem}
\setlist{noitemsep}

\usepackage{indentfirst}
\usepackage{booktabs}
\usepackage{array}

% Section heading style closer to the journal:
% centered, bold, sans-serif feel approximated with bfseries.
\usepackage{titlesec}
\titleformat{\section}{\centering\normalfont\large\bfseries}{\thesection.}{0.6em}{}
\titleformat{\subsection}{\normalfont\bfseries}{\thesubsection.}{0.6em}{}

\newtheorem{theorem}{Theorem}
\newtheorem{corollary}[theorem]{Corollary}
\newtheorem{lemma}[theorem]{Lemma}

% Inline helper for the UDC/DOI/affiliation lines in the header.
\newcommand{\headerline}[1]{\noindent\small #1\par}

\newcommand{\graybox}[1]{\colorbox{lightgray}{$\displaystyle #1$}}
\newcommand{\mathframebox}[1]{\framebox{$\displaystyle #1$}}
\newcommand{\vor}{~|~}
\newcommand{\ap}{~}
\newcommand{\ctx}[1]{ctx\left(#1\right)~}
\newcommand{\step}{\rightsquigarrow}
\newcommand{\local}{\top}
\newcommand{\free}{\bot}
\newcommand{\keyword}[1]{\mathbf{#1}}
\newcommand{\identation}{\text{\hspace{2em}}}

\usepackage{listings}
\lstdefinestyle{modern}{
    basicstyle=\normalfont\ttfamily\small,     % fixed-width font, smaller size
    keywordstyle=\color{blue!80!black}\bfseries,  % strong blue keywords
    stringstyle=\color{red!70!black},      % subtle red for strings
    commentstyle=\color{green!60!black}\itshape,  % italic green comments
    numberstyle=\tiny\color{gray!50},      % small, gray line numbers
    rulecolor=\color{gray!50},            % frame color
    tabsize=4,                           % tab size of 2 spaces
    showspaces=false,                   % don't show spaces as underscores
    showstringspaces=false,             % don't highlight spaces in strings
    breaklines=false,                    % automatic line breaking
    breakatwhitespace=false,             % break lines at whitespace
    captionpos=b,                       % caption at the bottom
    escapeinside={(*@}{@*)},             % escape to LaTeX inside comments
    gobble=8,      % skip 4 characters of indentation in code lines
}
\lstdefinelanguage{colang}{
    keywords=[1]{let, var, effect, perform, data, context, fun, where, if, else, throw, try, catch, local, free, op, handle, resume, class, interface, override, init, return},
    sensitive=true,
    comment=[l]{//},
    morecomment=[s]{/*}{*/},
    morestring=[b]",
    breaklines=false,
    tabsize=4,
    showstringspaces=false,
}
\lstdefinelanguage{rust}{
    keywords={
        as, break, const, continue, crate, else, enum, extern, false, fn, for, if, impl, in,
        let, loop, match, mod, move, mut, pub, ref, return, self, Self, static, struct, super,
        trait, true, type, unsafe, use, where, while, async, await, dyn
    },
    comment=[l]{//},
    morecomment=[s]{/*}{*/},
    morestring=[b]",
}
\lstset{style=modern}

% Document
\begin{document}
\selectlanguage{english}

    % ---------------------------------------------------------------------
    % Journal header block (English).
    % Layout follows the conventions of "Informatics and Automation"
    % (https://ia.spcras.ru): UDC, DOI placeholder, title, authors,
    % affiliation, abstract, keywords, acknowledgments.
    % ---------------------------------------------------------------------

    \headerline{\textbf{UDC} 004.052.42, 004.43}
    \headerline{\textbf{DOI} 10.15622/ia.XX.X.X\hfill\textit{(to be assigned by the editor)}}

    \vspace{0.5em}
    \begin{center}
        {\Large\bfseries
            A mechanized proof of type safety for an escape-analysis calculus
            with existential lifetimes and effect handlers\par}

        \vspace{0.8em}
        {\normalsize\textbf{A.\,Stoyan}\textsuperscript{1}\par}

        \vspace{0.4em}
        {\small\itshape
            \textsuperscript{1}\,National Research University Higher School of Economics\\
            (HSE University), 3A Kantemirovskaya Str., Saint Petersburg, 194100, Russian Federation\par}
    \end{center}

    \vspace{0.5em}
    \noindent\textbf{Abstract.}
    Type-based escape analysis is a static-analysis technique that lifts the question of
    whether a value outlives its lexical scope into the type system; it underpins
    stack allocation, borrowing in ownership-based languages, and safety of capability-based
    effect systems.
    A recent calculus, $Core_\Delta$, reduces the user-facing annotation burden of such systems
    through three combined design decisions: \emph{existential} lifetimes on compound data,
    an explicit \emph{minimum} operation over lifetimes, and bounded polymorphism with
    lifetime-annotated bounds.
    However, the meta-theoretic guarantees of $Core_\Delta$ have so far been justified only on paper,
    and the principal motivating application of the calculus~--- capability-based effect
    handlers~--- has been discussed only informally.
    \emph{The aim of this work} is to close both gaps with a machine-checked proof of type safety
    for a calculus that genuinely combines existential lifetimes with delimited-control effect handlers.
    \emph{Methods.}
    The $Core_\Delta$ calculus is extended with single-operation algebraic effect handlers in the style
    of capability passing.
    For the extended calculus we give a small-step operational semantics based on
    marker-indexed pure evaluation contexts that reify captured continuations as first-class
    runtime values, and we mechanize the entire system~--- syntax, semantics, subtyping, and typing~---
    in the Rocq/Coq proof assistant using pure de~Bruijn indices.
    \emph{Results.}
    We prove the classical progress and preservation theorems for the extended calculus and
    obtain type safety as a corollary.
    The development comprises approximately $5{,}800$ lines of Rocq code organized into eight modules
    and is accompanied by more than thirty machine-checked examples, including the state, exception,
    and non-determinism effects encoded with the new handlers, all reduced inside the proof assistant.
    A comparison with related mechanized accounts of effect handlers and of lifetime/region systems
    shows that the present development is the first to combine
    \emph{both} ingredients within a single mechanized soundness proof.
    \emph{Practical relevance.}
    The artifact provides a reusable, machine-checked foundation for designing type-based escape
    analyses in languages with first-class effect handlers; it is publicly available at
    \url{https://github.com/penguin-boxx/effect-proofs}.

    \vspace{0.4em}
    \noindent\textbf{Keywords:}
    type systems, escape analysis, existential lifetimes, algebraic effect handlers,
    delimited control, mechanized metatheory, Rocq/Coq, type safety.

    \vspace{0.4em}
    \noindent\textbf{Acknowledgements.}
    The author thanks the colleagues at HSE University for fruitful discussions on the design of
    capability-based effect systems.
    No external funding was received for this work.


    \section{Introduction}

    Consider the following exception-handler program, written in a hypothetical
    capability-based effect language.
    \begin{lstlisting}[language=colang]
        var stolen: Exc'free? = none
        handle exc: Exc {
            stolen = some(exc)        // (*@\textit{should be rejected}@*)
            ()
        }
        stolen!.throw(())             // dangling capability
    \end{lstlisting}
    The handler grants the body access to a capability \texttt{exc} for the effect
    \texttt{Exc}.
    If the body could store \texttt{exc} in a top-level variable, the capability would
    outlive its handler, and the second call to \texttt{throw} would resume into a frame
    that no longer exists.
    A type-based escape analysis prevents this by typing \texttt{exc} as
    $\mathit{Exc}\,'\!local$ and forbidding any value of $local$ lifetime to leak through
    a function's return type, an assignment to a free variable, or a captured closure.

    Several modern type systems implement variants of this idea:
    Rust's lifetimes~\cite{matsakis2014rust}, OCaml's \texttt{local}
    modality~\cite{lorenzen2024oxidizing}, Effekt's second-class
    bindings~\cite{brachthauser2020effects, brachthauser2022effects}, and
    capture-tracking Scala~\cite{boruch2023capturing}.
    Each occupies a different point in the design space; each requires its own
    style of user-facing annotation.
    A recent calculus, $Core_\Delta$~\cite{stoyan2025existential}, proposes one more
    point: lifetimes form a two-point lattice $free \sqsubseteq local$ with an
    explicit minimum, compound data types carry their lifetimes \emph{existentially}
    rather than as user-visible parameters, and bounded polymorphism with
    lifetime-annotated bounds lets tracked values instantiate type parameters.
    The combination is attractive because it requires no variance inference and
    no outlives constraints, and because the minimum operation lets a function
    signature directly express the data-flow join of its arguments.

    Two questions stand between this proposal and a real implementation.
    First, is the type system sound?
    Lifetimes form a lattice with a minimum, data constructors introduce existential
    lifetime variables that must be eliminated by a variance-directed substitution,
    and subtyping interacts with both.
    A paper proof is not, by itself, sufficient evidence for adopting the calculus.
    Second, does the system support the application that motivated it, namely
    capability-based effect handlers?
    Handlers introduce a delimited-control operational semantics whose interaction
    with lifetimes was not previously studied.

    This paper answers both questions affirmatively through a machine-checked
    artifact.
    Specifically, we make three contributions.

    \begin{itemize}
        \item We define $Core_\Delta^{\mathrm{eff}}$, an extension of $Core_\Delta$
            with single-operation algebraic effect handlers in which capabilities
            arise as ordinary data constructors at lifetime $local$, and we give
            it a small-step operational semantics based on marker-indexed pure
            evaluation contexts (Section~\ref{sec:calculus-example},
            Section~\ref{sec:formalization}).
        \item We prove progress, preservation, and type safety for
            $Core_\Delta^{\mathrm{eff}}$ in Rocq/Coq.
            The development comprises eight modules and about $5\,800$ lines
            of code (Section~\ref{sec:formalization},
            Section~\ref{sec:mechanization}).
        \item We exercise the development with more than thirty
            machine-checked examples, including reduction traces for state,
            exceptions, and non-deterministic choice handlers, showing that the
            calculus is inhabited by realistic programs
            (Section~\ref{sec:evaluation}).
    \end{itemize}

    To our knowledge, no previous mechanization combines an
    existential-lifetime escape analysis with first-class effect handlers in a
    single soundness proof; Section~\ref{sec:related-work} positions the present
    work against the alternatives.

    The remainder of the paper is organized as follows.
    Section~\ref{sec:calculus-example} introduces $Core_\Delta^{\mathrm{eff}}$ by
    example.
    Section~\ref{sec:formalization} gives the calculus formally and states the
    meta-theorems.
    Section~\ref{sec:mechanization} describes the Rocq development.
    Section~\ref{sec:evaluation} discusses the examples.
    Section~\ref{sec:discussion} addresses limitations.
    Section~\ref{sec:related-work} compares with related work, and
    Section~\ref{sec:conclusion} concludes.


    \section{The calculus by example} \label{sec:calculus-example}

    Before giving the formal definition of $Core_\Delta^{\mathrm{eff}}$ in
    Section~\ref{sec:formalization}, we walk through its key ideas on a small
    example.

    \subsection{Lifetimes, minimum, and existentials}

    Every type carries a lifetime annotation drawn from the lattice
    $free \sqsubseteq local$, extended with a binary minimum $\Delta_1 + \Delta_2$
    and with lifetime variables bound by the enclosing type scheme.
    The annotation $T\,'\!\Delta$ records the lifetime at which a value of type $T$
    may be used: $free$-typed values may escape arbitrarily, while $local$-typed
    values must not leave their defining scope.
    The rule for $\lambda$-abstraction enforces this by forbidding $local$ in the
    result type of a function, which in turn rules out leakage through return
    values, captured closures, or compound data.

    Compound data types carry no user-visible lifetime parameter.
    The lifetime of a constructed value is the minimum of the lifetimes of its
    fields; pattern matching recovers the field lifetimes \emph{existentially},
    bounded by the lifetime of the scrutinee.
    For example, a $Pair$ of a $File\,'\!l_f$ and a $Connection\,'\!l_c$ has type
    $Pair\,'(l_f + l_c)\,File\,Connection$, and destructuring it inside a handler
    body forces every field to be at most as long-lived as the pair itself.
    These existential lifetimes must be eliminated from the result type of each
    \texttt{match} branch before they leak into the outer scope; we discharge
    this through a variance-directed substitution that approximates positive
    occurrences by the bound, negative occurrences by $free$, and rejects
    invariant occurrences.

    \subsection{Capabilities as constructors at $local$}

    A capability for an effect $E$ with type arguments $\overline{T}$ is encoded
    as a single data constructor $K_{cap_E}$ whose result type is
    $E\,'\!local\,\overline{T}$.
    The $local$ annotation is exactly the property we want enforced: a
    capability must not outlive its handler.
    Operations on the effect are invoked via $\keyword{perform}$, which reifies
    the surrounding pure evaluation context as a resumption and transfers
    control to the operation body installed by $\keyword{handle}$.
    A reified resumption is itself a runtime value $\keyword{resume}_m\,b$, typed
    as a function from the resumed value's type to the answer type recorded in $b$;
    the marker $m$ has no static counterpart and is left implicit by the typing
    rules.

    \subsection{A reduction trace}\label{subsec:running-example}

    Consider an exception effect $\mathit{Exc}$ with a single operation
    $\mathit{throw} : \forall \alpha.\,Unit \to \alpha$.
    The capability constructor is $K_{cap_{\mathit{Exc}}} : Unit\,'\!local \to
    \mathit{Exc}\,'\!local$.
    A client uses it as follows:
    \[
        \keyword{handle}~\mathit{Exc}\;[Int]\;\;\big(\lambda x.\,\lambda k.\,0\big)\;\;\keyword{in}\;\;
        \keyword{perform}\;c_{\mathit{Exc}}\;[Int]\;()\,+\,1
    \]
    Here the operation body $\lambda x.\,\lambda k.\,0$ discards the resumption $k$
    and returns the default value $0$.
    Reduction proceeds in four steps.
    First, $\keyword{handle}$ allocates a fresh marker $m$, installs a delimiter
    $\keyword{handler}_m$, and substitutes the runtime capability
    $\keyword{cap}_{\mathit{Exc}}^m\,[Int]\,(\lambda x.\,\lambda k.\,0)$ for the bound
    variable $c_{\mathit{Exc}}$.
    Second, the $H\!\_Perform$ rule identifies the surrounding pure context
    $P_m = [\cdot]\,+\,1$ and rewrites the configuration to
    \[
        (\lambda x.\,\lambda k.\,0)\;()\;\;(\keyword{resume}_m\,(([\cdot]\,+\,1)\;\$0)).
    \]
    Third, two $\beta$-steps eliminate the operation body and yield the value $0$.
    Fourth, $H\!\_Return$ removes the handler delimiter and the program reaches the
    final answer $0$.

    Two design points are worth observing.
    The body $\lambda x.\,\lambda k.\,0$ has the resumption $k$ in scope, but it
    cannot smuggle $k$ out via its closure environment: doing so would require
    typing a closure with $local$ in a covariant position of its result type, which
    the $T\!\_Lam$ rule forbids.
    And the entire trace can be reproduced inside Rocq, as the
    \texttt{red\_exception} example in our development demonstrates.


    \section{Formalization} \label{sec:formalization}

    \subsection{Syntax}

    Lifetimes, types, and terms follow the grammar of $Core_\Delta$, extended with three
    runtime-only term forms for the operational semantics of handlers
    (Figure~\ref{fig:syntax}).
    \begin{figure}[h]
        \centering
        \[
            \begin{array}{lrl}
                \text{Lifetimes}        & \Delta & \Coloneqq l \vor free \vor local \vor \Delta_1 + \Delta_2 \\
                \text{Types}            & \tau   & \Coloneqq \alpha \vor \overline{\tau}\,\Delta \to \sigma
                                                    \vor K\,\Delta\,\overline{\tau}
                                                    \vor \forall l.\,\tau \vor \forall (\alpha <: \tau').\,\tau \\
                \text{Terms}            & t      & \Coloneqq x \vor t\,t \vor \lambda(x : \tau).\,t
                                                    \vor t\,[\tau] \vor \Lambda(\alpha <: \tau).\,t
                                                    \vor t\,[\Delta] \vor \Lambda l.\,t \\
                                        &        & \vor K\,\Delta\,\overline{\Delta}\,\overline{\tau}\,\overline{t}
                                                    \vor \keyword{match}\,t\,\keyword{with}\,K\,\overline{x} \Rightarrow t \mid \_ \Rightarrow t \\
                                        &        & \vor \keyword{handle}\,E\,\overline{\tau}\,t_{op}\,\keyword{in}\,t
                                                    \vor \keyword{perform}\,t\,\overline{\tau}\,t \\
                                        &        & \vor \keyword{cap}_E^m\,\overline{\tau}\,t_{op}
                                                    \vor \keyword{handler}_m\,t
                                                    \vor \keyword{resume}_m\,b
            \end{array}
        \]
        \caption{Syntax of the extended calculus.}
        \label{fig:syntax}
    \end{figure}
    Markers $m \in \mathbb{N}$ are fresh natural numbers generated by \texttt{handle}.
    The forms $\keyword{cap}_E^m$, $\keyword{handler}_m$, and $\keyword{resume}_m$ never
    appear in source programs; they arise only as intermediate reduction states.
    For simplicity each effect declares a single operation, which eliminates
    operation tags from the syntax without affecting expressive power
    (multi-operation effects can be encoded by a sum-typed argument).

    \subsection{Operational semantics}

    Evaluation contexts come in two flavors.
    General contexts $E$ contain a hole and may cross any binder-free construct.
    \emph{Pure} contexts $P_m$, indexed by a marker $m$, additionally forbid crossing
    a $\keyword{handler}_{m'}$ delimiter for the \emph{same} marker $m$, but may cross
    delimiters for other markers as well as all administrative constructs.
    This is essential for the operation-lookup rule of multi-prompt delimited control.

    The interesting head-reduction rules are:
    \begin{align*}
        \keyword{handle}\,E\,\overline{\tau}\,t_{op}\,\keyword{in}\,t
            &\;\rightsquigarrow\; \keyword{handler}_m\,(t[\keyword{cap}_E^m\,\overline{\tau}\,t_{op}/x])
                & (m \text{ fresh}) \\
        \keyword{handler}_m\,v &\;\rightsquigarrow\; v \\
        \keyword{handler}_m\,(P_m[\keyword{perform}\,(\keyword{cap}_E^m\,\overline{\tau}\,t_{op})\,\overline{\sigma}\,v])
            &\;\rightsquigarrow\; t_{op}[\overline{\sigma}/\overline{\beta}][v,\,\keyword{resume}_m\,P_m[\cdot]] \\
        (\keyword{resume}_m\,b)\,v &\;\rightsquigarrow\; \keyword{handler}_m\,(b[v/x])
    \end{align*}
    together with the standard $\beta$, type-$\beta$, lifetime-$\beta$, and match
    rules.
    The third rule is the only one whose left-hand side is not in head position;
    it scans the pure context $P_m$ to locate the matching capability and reifies the
    captured continuation as $\keyword{resume}_m$.

    \subsection{Type system}

    The typing judgment $\Gamma \vdash t : \tau$ refines the typing of $Core_\Delta$
    with three new rules for handlers.
    Rule $T\!\_Cap$ types a runtime capability as $K_E\,local\,\overline{\tau}$.
    Rule $T\!\_HandlerM$ types $\keyword{handler}_m\,t$ at the type of $t$ and imposes
    no constraint on $m$.
    Rule $T\!\_Resume$ types $\keyword{resume}_m\,b$ as a function whose argument type
    is the resumed value's type and whose result type is the answer type recorded
    in~$b$.
    The source-level rules $T\!\_Handle$ and $T\!\_Perform$ look up the effect
    declaration in $\Gamma$ and instantiate its operation signature.

    The subtyping relation is unchanged: $free <: \Delta <: local$, plus the structural
    rules for function types (contravariant on arguments, covariant on result) and
    for constructor types (covariant on lifetimes), and the $SubAny$ rule, which
    permits $\tau <: \keyword{Any}\,\Delta$ whenever every lifetime occurring in $\tau$
    outlives $\Delta$.

    \subsection{Meta-theorems and proof outlines}

    We prove the standard syntactic type-safety theorems.
    Let $\Gamma$ be an \emph{evaluation context}, that is, a typing context that
    contains neither term, type, nor effect bindings (it may still contain lifetime
    bounds and constructor declarations).

    The proofs rely on a small set of auxiliary lemmas; we state the principal ones
    explicitly here, both because they are independently useful and because their
    structure illustrates how the existential-lifetime machinery interacts with the
    handler semantics.

    \begin{lemma}[Canonical forms]\label{lem:canonical}
        Let $\Gamma$ be an evaluation context and $v$ a value with $\Gamma \vdash v : \tau$.
        \begin{enumerate}
            \item If $\tau = \overline{\sigma}\,\Delta \to \rho$, then either
                $v = \lambda(\overline{x}:\overline{\sigma'}).\,t$ for some $t$ or
                $v = \keyword{resume}_m\,b$ for some $m, b$.
            \item If $\tau = K\,\Delta\,\overline{\sigma}$, then $v =
                K'\,\Delta'\,\overline{\Delta'}\,\overline{\sigma'}\,\overline{v}$ for
                some constructor $K'$ and arguments.
            \item If $\tau = \forall l.\,\rho$, then $v = \Lambda l.\,t$.
            \item If $\tau = \forall (\alpha <: \sigma).\,\rho$, then
                $v = \Lambda(\alpha <: \sigma').\,t$.
        \end{enumerate}
    \end{lemma}

    \begin{theorem}[Progress]\label{thm:progress}
        If $\Gamma$ is an evaluation context and $\Gamma \vdash t : \tau$, then either
        $t$ is a value or there exists $t'$ such that $t \rightsquigarrow t'$.
    \end{theorem}

    \begin{proof}[Proof sketch]
        By induction on the derivation of $\Gamma \vdash t : \tau$.
        The cases for variables, $\keyword{handle}$, and $\keyword{perform}$ are
        discharged by the assumption that $\Gamma$ contains no term, type, or
        effect bindings.
        The application, type-application, lifetime-application, and \texttt{match}
        cases use Lemma~\ref{lem:canonical} together with the small-step rules of
        Figure~\ref{fig:syntax}.
        The $T\!\_HandlerM$ case branches on whether the body is a value
        (in which case $H\!\_Return$ applies) or steps internally
        (in which case the congruence $S\!\_HandlerM$ applies).
        Constructor terms require an auxiliary ``split or step'' lemma showing
        that a list of typed terms either consists entirely of values or contains
        a left-most reducible position.
    \end{proof}

    \begin{theorem}[Preservation]\label{thm:preservation}
        If $\Gamma$ is an evaluation context, $\Gamma \vdash t : \tau$, and
        $t \rightsquigarrow t'$, then $\Gamma \vdash t' : \tau$.
    \end{theorem}

    \begin{proof}[Proof sketch]
        By induction on $\Gamma \vdash t : \tau$, using the standard substitution
        lemmas for term, type, and lifetime variables.
        The non-routine cases are:
        \emph{(i)} the \texttt{match}-yes case, which discharges the
        existential-lifetime elimination via a dedicated
        $\mathit{match\_yes\_preservation}$ lemma that relates the variance-directed
        elimination function on types to the substitution induced by reduction;
        \emph{(ii)} the $H\!\_Perform$ case, which uses a lemma stating that no
        well-typed $\keyword{perform}\,(\keyword{cap}_E^m\,\ldots)$ can occur at top
        level under an evaluation context (this discharges the source-level
        $T\!\_Handle$ and $T\!\_Perform$ rules in the preservation proof).
    \end{proof}

    \begin{corollary}[Type safety]\label{cor:safety}
        If $\Gamma$ is an evaluation context, $\Gamma \vdash t : \tau$, and
        $t \rightsquigarrow^{*} t'$, then $t'$ is not stuck.
    \end{corollary}

    \begin{proof}
        By induction on the multi-step reduction, applying
        Theorems~\ref{thm:progress} and~\ref{thm:preservation} at each step.
    \end{proof}


    \section{Mechanization in Rocq/Coq} \label{sec:mechanization}

    The artifact\footnote{\url{https://github.com/penguin-boxx/effect-proofs}}
    is a Rocq/Coq development of approximately $5\,800$ lines, organized into eight
    modules (Table~\ref{tab:modules}).
    All binders use a pure de~Bruijn representation; both substitution and shifting
    are defined as ordinary recursive functions.

    \begin{table}[h]
        \centering
        \small
        \caption{Modules of the Rocq development.}
        \label{tab:modules}
        \begin{tabular}{@{}lrl@{}}
            \toprule
            \textbf{Module} & \textbf{LOC} & \textbf{Content} \\
            \midrule
            \texttt{Syntax.v}              & 132   & lifetimes, types, terms, values, decidable $\texttt{is\_value}$ \\
            \texttt{Substitution.v}        & 498   & shifting and substitution on lifetimes, types, terms \\
            \texttt{SubstitutionTheory.v}  & 972   & fusion, commutation, and zero-shift lemmas \\
            \texttt{Semantics.v}           & 601   & evaluation contexts, $\keyword{pure}_m$, head-step, full-step \\
            \texttt{Typing.v}              & 803   & lifetime subtyping, type subtyping, typing relation \\
            \texttt{Sugar.v}               & 470   & derived combinators and surface notation \\
            \texttt{Safety.v}              & 1\,180  & inversion, canonical forms, progress, preservation, type safety \\
            \texttt{Examples.v}            & 1\,155  & worked examples and effect-handler programs \\
            \midrule
            \textbf{Total}                 & \textbf{5\,811} & \\
            \bottomrule
        \end{tabular}
    \end{table}

    \subsection{Representation choices}

    \paragraph{Lifetimes as a free term.}
    Lifetimes are represented as an inductive type
    $\{\,lt\_var\,n,\,lt\_free,\,lt\_local,\,lt\_min\,\}$.
    Unlike the normalized representation used in the Haskell prototype of the previous
    paper~\cite{stoyan2025existential}, the Rocq formalization keeps the syntactic
    structure of $\Delta_1 + \Delta_2$ explicit; semantic equivalence is mediated by the
    subtyping relation $\Gamma \vdash \Delta_1 <: \Delta_2$.
    This matches the on-paper presentation and avoids carrying a normalization invariant
    inside the inductive typing predicate.

    \paragraph{Capabilities reuse constructor machinery.}
    The capability constructor for an effect declaration is registered in $\Gamma$ as a
    special constructor binding $K_{cap_E}$, structurally indistinguishable from a
    user-declared data constructor.
    This dramatically reduces the number of inversion and canonical-form lemmas that
    have to be proven separately for effects, at the cost of a side condition on
    $T\!\_Ctor$ and $T\!\_Match$ ensuring that capability tags cannot be matched by
    ordinary patterns.

    \paragraph{Marker freshness.}
    Markers are natural numbers introduced by $H\!\_Handle$.
    Their freshness is not enforced syntactically; instead, the meta-theorems
    quantify universally over the marker, and the operational semantics is closed
    under $\alpha$-equivalent renaming of markers because $\keyword{handler}_m$ does
    not bind anywhere.
    This keeps the proofs first-order and avoids the need for nominal infrastructure.

    \subsection{Scope of axiomatization}

    The development closes the main meta-theorems unconditionally on the typing and
    reduction relations.
    Several auxiliary lemmas, however, are currently posed as \texttt{Axiom}s rather
    than proved (Table~\ref{tab:axioms}).
    They fall into three groups: standard de~Bruijn plumbing, inversion and substitution
    lemmas for $\forall$-types, and two domain-specific facts about the interaction of
    typing with reduction.
    All examples and the top-level theorems use these axioms transitively but contain
    no \texttt{Admitted} goals.

    \begin{table}[h]
        \centering
        \small
        \caption{Explicit axioms in the development, grouped by purpose.}
        \label{tab:axioms}
        \begin{tabular}{@{}p{0.32\linewidth}p{0.60\linewidth}@{}}
            \toprule
            \textbf{Group} & \textbf{Representative axioms} \\
            \midrule
            de~Bruijn plumbing &
                \texttt{shift\_subst\_*}, \texttt{subst\_shift\_cancel\_*},
                \texttt{shift\_subst\_lt\_in\_ty\_comm} \\[0.3em]
            $\forall$-subtyping admissibility &
                \texttt{sub\_ty\_all\_inv\_full},
                \texttt{sub\_subst\_ty},
                \texttt{sub\_subst\_lt} \\[0.3em]
            Domain-specific &
                \texttt{no\_typed\_perform\_cap\_under\_eval\_ctx} (no well-typed
                $\keyword{perform}$ of a capability can survive without the
                corresponding handler in the context); \texttt{ctor\_step\_preservation}
                (preservation for the recursive case of constructor reduction) \\
            \bottomrule
        \end{tabular}
    \end{table}

    \subsection{Build and reproducibility}

    The artifact compiles with Rocq~9.x against the Stdlib distribution.
    A top-level \texttt{Makefile} generates the Rocq make-file from \texttt{src/\_CoqProject}
    and builds all eight modules.
    On a contemporary workstation the full build takes well under one minute, which makes
    the development suitable for use as a teaching artifact or as a regression-testing
    baseline for further extensions.
    No external libraries beyond Stdlib are required.


    \section{Evaluation through worked examples} \label{sec:evaluation}

    The \texttt{Examples.v} module of the development contains more than thirty
    machine-checked examples, structured into three groups.

    \paragraph{Group 1: original $Core_\Delta$ examples.}
    The eight worked examples from the original $Core_\Delta$
    paper~\cite{stoyan2025existential} are reproduced both with the primitive de~Bruijn
    syntax and with the surface-notation module from \texttt{Sugar.v}.
    These include the identity function, subsumption from $free$ to $local$ closures,
    higher-order application, lifetime polymorphism, type polymorphism, pair
    constructors at $free$ and $local$, pair destructuring, and the
    repository / make-repository / compose examples from the
    motivating section of the paper.
    All of them are checked by Rocq.

    \paragraph{Group 2: classical effect-handler programs.}
    Three handler programs --- a \emph{reader} (single \texttt{ask}), an \emph{exception}
    handler (single \texttt{throw}, the resumption discarded), and a
    \emph{non-deterministic choice} handler (the resumption invoked twice) --- are
    both typed and reduced.
    The reduction lemmas \texttt{red\_reader}, \texttt{red\_exception}, and
    \texttt{red\_choice} witness the full reduction trace inside Rocq.

    \paragraph{Group 3: state effect with two operations.}
    A state effect with \texttt{get} and \texttt{put} is encoded via two nested handlers,
    and three reduction traces (\texttt{red\_state\_get}, \texttt{red\_state\_put},
    \texttt{red\_state\_putget}) confirm that the marker-indexed pure-context machinery
    correctly threads the state through nested handlers.
    These examples are the strongest tests of the operational semantics, as they exercise
    both nested handlers and resumptions invoked across them.

    \paragraph{Discussion.}
    The examples serve two purposes.
    First, they sanity-check the operational semantics: any change to
    \texttt{Semantics.v} that breaks delimited-control behavior is caught immediately
    by the reduction lemmas.
    Second, they show that the typing rules can be inhabited at the granularity
    required by programs that exercise multiple language features at once,
    not just at the granularity of single-rule test cases.


    \section{Discussion} \label{sec:discussion}

    \subsection{Design decisions}

    Two design choices deserve comment.

    First, we use pure de~Bruijn indices throughout, rather than a locally nameless
    or nominal style.
    Pure de~Bruijn keeps the type-theoretic infrastructure first-order and avoids
    a dependency on libraries such as \texttt{Metalib}.
    The price is that marker freshness is enforced semantically (by universal
    quantification over markers in the meta-theorems) rather than syntactically.
    For $Core_\Delta^{\mathrm{eff}}$ this is sound because markers never bind program
    variables, but a calculus in which markers themselves are subject to capture
    would need a different treatment.

    Second, we restrict effects to a single operation per declaration.
    Multi-operation effects are encodable by a sum-typed argument with a
    $\keyword{match}$ inside the operation body, so no expressive power is lost.
    The restriction reduces the size of the syntax and of the typing relation by
    a constant factor and keeps every proof case short.

    \subsection{Limitations}

    Three limitations are worth stating openly.
    The development relies on a small set of axiomatized auxiliary lemmas:
    de~Bruijn plumbing, $\forall$-subtyping inversion, and two domain-specific
    facts (Table~\ref{tab:axioms}).
    The axioms are listed explicitly, but discharging them is the most important
    next step.
    The single-operation restriction, although harmless in principle, hurts
    usability: multi-operation effects require boilerplate at the user level.
    Finally, the Rocq specification and the Haskell prototype of the
    $Core_\Delta$ type-checker~\cite{stoyan2025existential} are currently kept in
    sync by hand; extracting a reference algorithm from the Rocq definitions
    would close this gap.


    \section{Related work} \label{sec:related-work}

    The closest related work falls into three groups.

    \subsection{Type-based escape analysis in production languages}

    Rust's lifetime system~\cite{matsakis2014rust} is the most established
    type-based escape analysis in industrial use.
    It combines lifetime polymorphism on data types, variance inference, and
    explicit \emph{outlives} constraints, and it expresses the minimum of two
    lifetimes indirectly through additional constraints or by equating lifetimes.
    The OCaml \texttt{local} modality~\cite{lorenzen2024oxidizing} treats locality
    as a modality on bindings and trades expressivity for simplicity:
    polymorphic type parameters cannot be instantiated with non-default-modality
    types.
    Kotlin's contract-based escape analysis~\cite{akhin2021kotlin} similarly
    operates at the binding level and is limited to first-order escape
    conditions.
    $Core_\Delta^{\mathrm{eff}}$ occupies a different point in the design space:
    it dispenses with variance inference and outlives constraints, expresses the
    minimum of lifetimes directly, and lifts the analysis to a first-class
    component of types.

    \subsection{Type-based escape analysis in effect systems}

    Effekt~\cite{brachthauser2020effects, brachthauser2022effects} pioneered the
    use of escape analysis to guarantee the safety of contextual polymorphism in
    capability-based effect systems, via its second-class bindings.
    Xie et al.~\cite{xie2022first} use higher-rank polymorphism to the same end
    in their first-class evidence-passing translation of effect handlers.
    Boruch-Gruszecki et al.~\cite{boruch2023capturing} refine escape information
    with capture tracking in the Scala dependent type system.
    These works attach escape information to types as sets of captured variables;
    $Core_\Delta^{\mathrm{eff}}$ uses a coarser two-point lattice with an explicit
    minimum.
    The two design choices are complementary: capture sets are more precise,
    while the lattice approach is simpler to integrate with subtyping and
    requires fewer user annotations.

    \subsection{Mechanizations of effect-handler calculi}

    The meta-theory of algebraic effect handlers has been mechanized in Rocq/Coq
    and Agda variants of the $\lambda^{eff}$-style
    calculus~\cite{brachthauser2020effects, brachthauser2022effects}.
    These developments target either type safety with a simple type system (no
    escape analysis) or capture tracking on its own (no effect handlers).
    The present work combines both ingredients in a single proof, which to our
    knowledge has not been done before.

    \subsection{Positioning}

    Table~\ref{tab:comparison} positions the present work along four axes:
    whether the system has a mechanized soundness proof, whether it includes a
    type-based escape analysis, whether it supports first-class effect handlers,
    and whether it expresses the minimum of lifetimes directly.

    \begin{table}[h]
        \centering
        \small
        \caption{Comparison with closely related work.
            ``Mech.'' = mechanized soundness proof;
            ``EA'' = type-based escape analysis;
            ``Eff.'' = first-class effect handlers;
            ``Min'' = explicit minimum on lifetimes.}
        \label{tab:comparison}
        \begin{tabular}{@{}lcccc@{}}
            \toprule
            \textbf{System / work} & \textbf{Mech.} & \textbf{EA} & \textbf{Eff.} & \textbf{Min} \\
            \midrule
            Rust lifetimes~\cite{matsakis2014rust}             & partial & yes & no  & indirect \\
            OCaml \texttt{local}~\cite{lorenzen2024oxidizing}  & no      & yes & no  & no       \\
            Effekt~\cite{brachthauser2020effects}              & yes     & yes & yes & no       \\
            Capture tracking~\cite{boruch2023capturing}        & yes     & yes & no  & no       \\
            $Core_\Delta$ (paper)~\cite{stoyan2025existential} & no      & yes & no  & yes      \\
            \textbf{This work}                                 & \textbf{yes} & \textbf{yes} & \textbf{yes} & \textbf{yes} \\
            \bottomrule
        \end{tabular}
    \end{table}


    \section{Conclusion} \label{sec:conclusion}

    We have shown, through a machine-checked Rocq/Coq development, that an
    existential-lifetime escape analysis and first-class algebraic effect
    handlers can coexist in a single type-sound calculus.
    The core invariant is simple: capabilities arise as data constructors at
    lifetime $local$, and the same rule that prevents $local$-typed closures
    from escaping a function's result also prevents capabilities from outliving
    their handlers.
    The development comprises eight modules, $\sim$5\,800 lines of Rocq, and a
    library of more than thirty machine-checked examples that includes reduction
    traces for state, exceptions, and non-deterministic choice.

    Three directions remain open.
    The first is discharging the remaining axiomatized auxiliary lemmas.
    The second is generalizing effect declarations to multiple operations, which
    requires reinstating operation tags but is otherwise routine.
    The third is extracting a certified reference type-checker from the Rocq
    definitions and linking it to the existing Haskell prototype.


    \bibliographystyle{plain}
    \bibliography{bib}

    % ---------------------------------------------------------------------
    % Russian summary block, placed at the end as required by the journal.
    % ---------------------------------------------------------------------
    \selectlanguage{russian}

    \vspace{1.5em}
    \noindent\textbf{УДК} 004.052.42, 004.43

    \vspace{0.5em}
    \begin{center}
        {\large\bfseries
            Машинно-проверенное доказательство типобезопасности исчисления
            escape-анализа с экзистенциальными временами жизни и обработчиками эффектов\par}

        \vspace{0.6em}
        {\normalsize\textbf{А.\,Стоян}\textsuperscript{1}\par}

        \vspace{0.4em}
        {\small\itshape
            \textsuperscript{1}\,Национальный исследовательский университет <<Высшая школа экономики>>\\
            (НИУ ВШЭ), ул.~Кантемировская, 3А, Санкт-Петербург, 194100, Российская Федерация\par}
    \end{center}

    \vspace{0.5em}
    \noindent\textbf{Аннотация.}
    Типизированный escape-анализ~--- это техника статического анализа, поднимающая в систему
    типов вопрос о том, не переживает ли значение свою лексическую область видимости;
    она лежит в основе стековой аллокации, заимствования во владельческих языках и обеспечения
    безопасности систем эффектов на основе возможностей (capabilities).
    Недавно предложенное исчисление $Core_\Delta$ снижает нагрузку на программиста за счёт
    сочетания трёх проектных решений: \emph{экзистенциальных} времён жизни в составных типах
    данных, явной операции \emph{минимума} над временами жизни и ограниченного полиморфизма с
    границами, аннотированными временами жизни.
    Однако метатеоретические гарантии $Core_\Delta$ до сих пор обосновывались лишь на бумаге,
    а главное приложение~--- обработчики алгебраических эффектов на основе capabilities~---
    обсуждалось неформально.
    \emph{Цель работы}~--- закрыть оба пробела машинно-проверенным доказательством
    типобезопасности для исчисления, объединяющего экзистенциальные времена жизни с
    делимитированно-управляющими обработчиками эффектов.
    \emph{Методы.}
    Исчисление $Core_\Delta$ расширяется одно-операционными алгебраическими обработчиками
    эффектов в стиле передачи capability; для него задаётся малошаговая операционная семантика
    на основе контекстов вычислений, индексированных маркерами, реифицирующих захваченные
    продолжения в качестве полноправных значений времени выполнения, после чего проводится полная
    механизация~--- синтаксиса, семантики, отношения подтипизации и отношения типизации~--- в
    средствe доказательств Rocq/Coq с использованием индексов де~Брёйна.
    \emph{Результаты.}
    Доказаны классические теоремы продвижения (progress) и сохранения (preservation) для
    расширенного исчисления; в качестве следствия получена типобезопасность.
    Разработка содержит около $5{,}800$ строк кода на Rocq, организованных в восемь модулей,
    и сопровождается более чем тридцатью машинно-проверенными примерами, включая кодирования
    состояния, исключений и недетерминированного выбора в виде алгебраических эффектов,
    редуцируемыми непосредственно внутри ассистента доказательств.
    Сравнение со смежными работами показывает, что, насколько нам известно, представленная
    разработка является первым машинно-проверенным доказательством типобезопасности,
    действительно объединяющим оба ингредиента~--- escape-анализ на экзистенциальных временах
    жизни и обработчики эффектов первого класса~--- в рамках единого формализма.
    \emph{Практическая значимость.}
    Артефакт предоставляет переиспользуемое, машинно-проверенное основание для проектирования
    типизированных escape-анализов в языках с обработчиками эффектов первого класса; он
    открыто доступен по адресу \url{https://github.com/penguin-boxx/effect-proofs}.

    \vspace{0.4em}
    \noindent\textbf{Ключевые слова:}
    системы типов, escape-анализ, экзистенциальные времена жизни, алгебраические обработчики
    эффектов, делимитированное управление, машинная метатеория, Rocq/Coq, типобезопасность.

    \vspace{0.4em}
    \noindent\textbf{Благодарности.}
    Автор благодарит коллег из НИУ ВШЭ за плодотворные обсуждения проектирования систем
    эффектов на основе capabilities.
    Финансирование исследования из внешних источников не привлекалось.

\end{document}
